\textbf{结论名称：} 祖暅原理（幂势既同则积不容异）

\textbf{适用条件：} 两个几何体夹在两平行平面之间，任一平行于这两平面的平面截得的截面面积处处相等，且两几何体的高度相同。

\textbf{变量说明：} 设两几何体 $A,B$ 被夹在两平行平面间，高均为 $H$。以平行底面的平面截几何体，截面距底面高度为 $h\ (0\le h\le H)$，截面面积分别为 $S_A(h),S_B(h)$；两几何体体积记为 $V_A,V_B$。

\textbf{核心公式：} 
\[
\boxed{\forall\, h\in[0,H],\ S_A(h)=S_B(h) \ \Longrightarrow\ V_A=V_B}
\]

\textbf{使用场景：} 通过构造截面面积处处相等的已知几何体，间接求不规则几何体的体积；用于推导柱、锥、球等的体积公式。

\textbf{简单例子：} 
(1) 底面积为 $S$、高为 $h$ 的任意斜棱柱，与同底等高的直棱柱在任意高度处的截面均为原底面，面积恒为 $S$。由祖暅原理知二者体积相等，故斜棱柱体积 $V=Sh$。
(2) 求半径为 $R$ 的球体积：用高为 $R$ 的半球与底面半径为 $R$、高为 $R$ 的圆柱挖去同底等高圆锥后的几何体比较，同一高度处截面面积恒相等，得半球体积 $V_{\text{半球}}=\frac{2}{3}\pi R^3$，从而球体积 $V_{\text{球}}=\frac{4}{3}\pi R^3$。